\documentclass[12pt, oneside]{article}

\usepackage{xparse}

\usepackage{geometry}
\geometry{letterpaper}
\usepackage[parfill]{parskip}
\usepackage{float}

% Math packages
\usepackage{amssymb}
\usepackage{amsmath}
\usepackage{amsthm}
\usepackage{mathalpha}

% Tables
\usepackage{booktabs}
\usepackage{afterpage}

% Hyperlinks
\usepackage{hyperref}

% References
\numberwithin{equation}{section}
\usepackage[numbers]{natbib}

% TikZ (optional for later)
\usepackage{tikz}
\usetikzlibrary{shapes,arrows,chains}

% Theorem environments
\theoremstyle{definition}
\newtheorem{definition}{Definition}[section]
\newtheorem{proposition}{Proposition}[section]
\newtheorem{axiom}{Axiom}[section]

\theoremstyle{plain}
\newtheorem{theorem}{Theorem}[section]
\newtheorem{lemma}[theorem]{Lemma}
\newtheorem{corollary}[theorem]{Corollary}

\theoremstyle{remark}
\newtheorem*{remark}{Remark}

\title{Affine Constraint Dynamics and Density-One Descent in the Collatz Map}
\author{Wayne Brassem}

\begin{document}

% Custom commands
\NewDocumentCommand{\setN}{}{\mathbb{N}}
\NewDocumentCommand{\setZ}{}{\mathbb{Z}}
\NewDocumentCommand{\setQ}{}{\mathbb{Q}}

% Document commands which provide hyperlinks to OEIS sequences referenced herein
\NewDocumentCommand{\OEIS}{m}{\href{https://oeis.org/#1}{#1}}

\maketitle

\begin{abstract}

We study the Collatz map via the fully accelerated transformation
\[
f(x)=\frac{3x+1}{2^{v_2(3x+1)}}.
\]

Finite trajectories are encoded by division words, each inducing an affine map
whose admissibility imposes congruence constraints on initial values. This
yields a symbolic--affine framework in which admissible words correspond to
compatible residue classes determined by inverse affine lifting.

We analyze the growth of admissible words through their associated division
counts, showing that the induced affine constraints refine at an exponential
rate governed by $2^{D(\omega)}$, while symbolic expansion remains subcritical.
This imbalance forces a sparsification of admissible initial conditions.

As a consequence, the set of integers that fail to admit a finite descent prefix
has natural density zero. Density-one descent follows as a corollary of this
affine growth domination.

\end{abstract}

\newpage
\tableofcontents

% ============================================================
\newpage
\section{Introduction}

Below is sample excerpt from the old spine. It is not yet fully integrated into
the new structure, but it gives a sense of the motivation and context for the work and
to verify ToC, reference, formating, pagination, etc.

The $3x+1$ problem, also known as the Collatz conjecture~\cite{Collatz1937}, concerns
the behavior of the map
\[
T(x) =
\begin{cases}
\displaystyle
    \phantom{3x} \frac{x}{2} & \text {if } x \text{ is even}, \\[8pt]
\displaystyle
    3x+1,                    & \text{if $x$ is odd}.
\end{cases}
\]
on the positive integers. It asserts that every positive integer eventually reaches $1$
under iteration of $T$. Despite its elementary formulation, the conjecture has resisted
proof for decades~\cite{Lagarias2010,Terras1976}. Extensive computational verification
confirms convergence for very large ranges~\cite{OliveiraESilva2010}, but does not
illuminate the global structure of trajectories.

etc., etc....
% High-level motivation
% Shift from measure-theoretic to affine-symbolic viewpoint
% Summary of main result and structure

% ============================================================
\section{Symbolic Encoding of Collatz Trajectories}
\label{sec:division-counts}

\subsection{Division Words}

\begin{definition}[Division Word]
A division word of length $m \ge 0$ is a sequence
\[
\omega = (d_0, \dots, d_{m-1}),
\]
where each $d_i \in \mathbb{N}$ represents the number of factors of $2$
removed in the $i$-th iterate of the accelerated map
\[
f(x) = \frac{3x+1}{2^{v_2(3x+1)}}.
\]
\end{definition}

\begin{remark}
Since $f$ maps odd integers to odd integers, each $3f^i(x)+1$ is even,
so $d_i \ge 1$ for all $i$.
\end{remark}

\begin{definition}[Realization]
An integer $x \in \mathbb{N}$ realizes $\omega$ if
\[
d_i = v_2(3f^i(x)+1), \quad 0 \le i < m.
\]
\end{definition}

\subsection{Total Division Count}

\begin{definition}
The total division count of a word $\omega$ is
\[
D(\omega) := \sum_{i=0}^{m-1} d_i.
\]
\end{definition}

\begin{remark}
The quantity $D(\omega)$ measures the total dyadic contraction
associated with the symbolic path $\omega$.
\end{remark}

% ============================================================
\section{Affine Representation of Division Words}

\begin{definition}
Each division word $\omega$ of length $m$ induces a formal affine map
\[
F_\omega(x) = \frac{3^m x + c(\omega)}{2^{D(\omega)}},
\]
obtained by symbolic composition of the accelerated map.
\end{definition}

\begin{proposition}
The map $\omega \mapsto F_\omega$ defines a semigroup homomorphism
from concatenation of division words to composition of affine maps.
\end{proposition}

\subsection{Example}

Consider convergent segments with $m=4$ iterates of $f$.  
The minimal total division count for $m=4$ is
\[
A020914(4) = 7.
\]

One admissible division word attaining this minimum is
\[
\omega = (1,1,1,4),
\]
for which
\[
D(\omega) = 1+1+1+4 = 7.
\]

Composing the accelerated map symbolically yields
\[
f_\omega(x)=\frac{3^4 x + c(\omega)}{2^7}
           = \frac{81x + 65}{128},
\]
so the canonical affine constant for this word is
\[
c(\omega)=65.
\]

The integrality condition
\[
81x + 65 \equiv 0 \pmod{128}
\]
determines a unique residue class
\[
x \equiv 15 \pmod{128}.
\]
This congruence defines a constraint fiber with modulus $2^7$.

\medskip

\noindent
\textbf{Interpretation.}
The word $\omega = (1,1,1,4)$ means:

\begin{itemize}
    \item At each of the first three iterates of $f$, $3x+1$ contains exactly one factor of $2$.
    \item At the fourth iterate of $f$, $3x+1$ contains four factors of $2$.
\end{itemize}

Under the accelerated map
\[
f(x)=\frac{3x+1}{2^{v_2(3x+1)}},
\]
all of these divisions occur within a single iterate. Thus $d_3=4$ means that
the fourth iterate divides by $2^4$ in that step. The constant $c(\omega)$ is
fixed entirely by this symbolic structure and ensures that precisely the
integers in the residue class $x \equiv 15 \pmod{128}$ produce integer iterates
under the affine form.

\medskip

This example illustrates that the behavior of trajectories is governed
entirely by the symbolic data encoded in $\omega$ and its associated
affine map. The admissible initial values are precisely those satisfying
the induced congruence constraint, independent of any particular choice
of representative.

% ============================================================
\section{Affine Symbolic Admissibility Grammar}

\begin{remark}[Division Grammar as a Dynamical Encoding]
Recall from Section~\ref{sec:division-counts} that a \emph{division word}
$\omega=(d_0,\dots,d_{m-1})$ records the exact 2-adic valuation
$d_i = v_2(3 x_i + 1)$, occurring at each iterate of the fully accelerated
Collatz map.

Unlike a purely combinatorial encoding, division words retain essential
dynamical information. Each division word $\omega=(d_0,\dots,d_{m-1})$
determines the sequence of multiplicative factors of the linear coefficient
\[
r_i = \frac{3}{2^{d_i}},
\]
so that the contractive or expansive contribution of each iterate is
preserved in the linear coefficient (with the affine term handled separately
via $c(\omega)$). This structure enables quantitative control of admissible
trajectories via affine growth and to formulate sufficient conditions for
descent, even though individual trajectories are collapsed under the encoding.
\end{remark}

\subsection{Explicit Form of the Affine Constant}
\begin{lemma}[Explicit Form of the Affine Constant]
Let $\omega = (d_0,\dots,d_{m-1})$ be an admissible division word.  
Then the affine map
\[
F_\omega(x) = \frac{3^m x + c(\omega)}{2^{D(\omega)}}
\]
has the explicit form
\[
c(\omega)
=
\sum_{j=0}^{m-1}
3^{m-1-j}\cdot 2^{\sum_{i=0}^{j-1} d_i},
\]
where empty sums are defined as $0$.
\end{lemma}

\subsection{Injectivity of the Affine Encoding}
\begin{lemma}[Injectivity of the Affine Encoding at Fixed Length]
\label{lem:injectivity}
Let $\omega = (d_0,\dots,d_{m-1})$ and $\omega' = (d'_0,\dots,d'_{m-1})$
be admissible division words of the same length $m$. Let
\[
c(\omega)
=
\sum_{j=0}^{m-1}
3^{m-1-j}\cdot 2^{D_j},
\quad
D_j = \sum_{i=0}^{j-1} d_i,
\]
and define $c(\omega')$ analogously.

If $\omega \ne \omega'$, then
\[
c(\omega) \ne c(\omega').
\]

Equivalently, the map $\omega \mapsto c(\omega)$ is injective on admissible
division words of fixed length $m$.
\end{lemma}

\begin{proof}
Suppose $\omega \ne \omega'$, and let
\[
k = \min\{\, j : d_j \ne d'_j \,\}
\]
be the first index at which the two words differ. Then for all $j < k$,
we have $D_j = D'_j$, while $D_k \ne D'_k$.

Consider the difference
\[
\Delta := c(\omega) - c(\omega')
=
\sum_{j=0}^{m-1}
3^{m-1-j}\bigl(2^{D_j} - 2^{D'_j}\bigr).
\]
The terms with $j<k$ cancel, so
\[
\Delta
=
3^{m-1-k}\bigl(2^{D_k} - 2^{D'_k}\bigr)
+
\sum_{j=k+1}^{m-1}
3^{m-1-j}\bigl(2^{D_j} - 2^{D'_j}\bigr).
\]

Let $\delta_k = \min(D_k, D'_k)$. Then
\[
2^{D_k} - 2^{D'_k}
=
2^{\delta_k}(2^a - 2^b),
\]
where one of $a,b$ is zero and $2^a - 2^b$ is odd. Hence
\[
v_2\!\left(3^{m-1-k}\bigl(2^{D_k} - 2^{D'_k}\bigr)\right)
=
\delta_k,
\]
since $3^{m-1-k}$ is odd.

For $j>k$, both $D_j$ and $D'_j$ strictly exceed $\delta_k$, since the
prefix sums increase by positive integers. Therefore
\[
v_2\!\left(3^{m-1-j}\bigl(2^{D_j} - 2^{D'_j}\bigr)\right)
>
\delta_k
\quad \text{for all } j>k.
\]

It follows that $\Delta$ is a sum of one term with 2-adic valuation
exactly $\delta_k$ and remaining terms of strictly higher 2-adic
valuation. Hence
\[
v_2(\Delta) = \delta_k,
\]
and in particular $\Delta \ne 0$.

Therefore $c(\omega) \ne c(\omega')$, completing the proof.
\end{proof}

This implies distinct division words induce distinct affine constraints.

\subsection{Stepwise Realizability}

\begin{theorem}[Stepwise Realizability of Division Words]
\label{thm:stepwise-realizability}
Let $\omega = (d_0,\dots,d_{m-1})$ be a division word, and define $c(\omega)$ by
\[
c(\omega)
=
\sum_{j=0}^{m-1}
3^{m-1-j}\cdot 2^{\sum_{i=0}^{j-1} d_i}.
\]

Then the identity
\[
2^{D(\omega)} x_m = 3^m x_0 + c(\omega)
\]
holds if and only if the sequence $(x_0,\dots,x_m)$ satisfies the accelerated
Collatz relations
\[
x_{i+1} = \frac{3x_i+1}{2^{d_i}}, \quad 0 \le i < m,
\]
with each intermediate step integral.

In particular, the congruence
\[
3^m x_0 + c(\omega) \equiv 0 \pmod{2^{D(\omega)}}
\]
is equivalent to the existence of a trajectory realizing $\omega$.
\end{theorem}

\begin{remark}[Structural Interpretation]
Each term in $c(\omega)$ corresponds to the additive contribution of a single
``+1'' in the Collatz iteration, propagated forward by subsequent multiplications
by $3$ and backward-scaled by preceding divisions by $2$.
\end{remark}

\subsection{Admissibility of Division Words}

\begin{definition}[Admissibility]
\label{def:admissible-division-word}
A division word $\omega$ is admissible if the congruence
\[
3^m x + c(\omega) \equiv 0 \pmod{2^{D(\omega)}}
\]
has a solution.
\end{definition}
\begin{remark}
Admissibility is therefore an arithmetic condition on the word $\omega$
itself, expressed as a divisibility constraint on the associated affine
map. When this condition holds, the set of all integers $x$ for which
$F_\omega(x)$ is integral defines a unique congruence constraint modulo
$2^{D(\omega)}$, corresponding to the residue class induced by $\omega$.
\end{remark}

\subsection{Minimal Division Count}

\begin{lemma}[Minimal Division Count]
For any admissible division word $\omega=(d_0,\dots,d_{m-1})$ of length $m$,
\[
d_0 + \cdots + d_{m-1} \;\ge\; A020914(m),
\]
where
\[
A020914(m) = \lceil m \log_2(3) \rceil.
\]
Equality holds if and only if $\omega$ achieves the minimal possible total
division count among admissible words of length $m$.
\end{lemma}

\subsection{Convergent Division Word}

\begin{definition}[Convergent Division Word]
An admissible division word $\omega = (d_0,\dots,d_{m-1})$ is said to be
\emph{convergent} (in the linear sense) if
\[
\frac{3^m}{2^{d_0+\cdots+d_{m-1}}} < 1.
\]
\end{definition}

% ============================================================
\section{Numerator Growth and Counting Control}

\begin{proposition}
The number of convergent division words of length $m$ is given by \OEIS{A186009}.
This sequence therefore describes the combinatorial growth rate of convergent
division words at depth $m$, and hence of the associated residue classes.
\end{proposition}

\begin{remark}[Canonical Tabular Realization]
A canonical realization of this counting sequence via the horizontal-sum construction,
together with its derivation from the increment grammar induced by \OEIS{A022921}, is given
in Appendix~\ref{app:horizontal-sum-table}. The appendix serves as a structural reference:
the arguments in this section rely only on the grammar itself and not on any numerical
properties of the tabulation.
\end{remark}

\begin{lemma}[Controlled Growth of Horizontal Sums]
\label{lem:controlled-horizontal-growth}

Let a row have strictly positive entries with total sum $S$.

A single increment step produces a new row whose total sum is $S$.

If a double increment is applied immediately afterward, producing a row in which
the terminal entry equals the sum of the previous row and all other entries remain
strictly positive, then the total sum of the resulting row satisfies
\[
S_{\mathrm{new}} < 3S.
\]
\end{lemma}

\newpage
\begin{proof}
The total sum of the original row is $S$.

Under a double increment step, the resulting row consists of:
\begin{enumerate}
    \item (i) interior entries with total $S_{\mathrm{int}} < S$, and
    \item (ii) two terminal contributions, each equal to $S$.
\end{enumerate}

Hence the total sum satisfies
\[
S_{\mathrm{new}} = S_{\mathrm{int}} + 2S < 3S.
\]
\end{proof}

\begin{theorem}[Maximal Admissible Growth Rate]
\label{thm:maximal-growth}
Let $\widetilde A(i) = A186009(i+1)$ denote the number of convergent residue classes
of stopping time $i$, generated by the horizontal-sum construction associated with
the increment grammar of \OEIS{A022921}.  Let $B(i) := A020914(i)$ denote the
cumulative exponent of $2$ up to level $i$.  

% The per-step growth of $\widetilde A(i)$ is bounded by the maximal geometric growth
% rate of a $5$--cycle (see Appendix~\ref{app:local-growth-details}, Equation~\eqref{eq:rho-5}):
% \[
% \limsup_{i\to\infty} \widetilde A(i)^{1/i} \le \rho_5 < 3,
% \]
% and consequently
% \[
% \frac{\widetilde A(i)}{2^{B(i)}} \;\longrightarrow\; 0 \quad \text{as } i\to\infty.
% \]
Then there exists a constant $\rho < 3$ such that
\[
\widetilde A(i) \le C \rho^i,
\]
and consequently
\[
\frac{\widetilde A(i)}{2^{B(i)}} \to 0.
\]
\end{theorem}

\begin{proof}
The increment sequence \OEIS{A022921} induces a finite symbolic grammar on $\{s,d\}$,
corresponding to single and double increments in the horizontal-sum construction.

% Among all admissible sequences, the shortest potentially expansive configuration is the
% \emph{5-cycle}
% \[
% \{s,d,s,d,d\}.
% \]
% All admissible words of length $\le 4$ produce horizontal-sum growth factors with
% geometric mean $\le 1$. A complete enumeration of admissible local patterns and
% their growth factors is given in Appendix~\ref{app:local-growth-details}.

% Explicit computation of the net horizontal-sum growth factor for the 5-cycle gives
% \[
% \rho_5 := \bigl(\text{horizontal sum multiplier over the 5-cycle}\bigr)^{1/5} < 3.
% \]
A detailed analysis of the increment grammar (see Appendix~\ref{app:local-growth-details})
shows that all admissible increment sequences have asymptotic growth rate bounded by a
uniform constant $\rho < 3$.

We denote the optimal such bound by $\rho_5 < 3$.

Any admissible infinite increment sequence decomposes into concatenations of finite
blocks permitted by the grammar. As explained in Remark~\ref{rem:grammar-suppression},
all locally expansive blocks are separated by mandatory contractive buffering. Horizontal-sum
growth is submultiplicative under concatenation, so the asymptotic per-step geometric
growth is bounded above by the maximal geometric mean among admissible blocks, which is
attained by the $5$-cycle and denoted $\rho_5$. Hence
\[
\limsup_{i\to\infty} \widetilde A(i)^{1/i} \le \rho_5 < 3.
\]

From Lemma~\ref{lem:uniform-control}, the denominator satisfies
\[
3^i < 2^{B(i)} < 2\cdot 3^i.
\]
Consequently, there exists a constant $C>0$, depending only on initial conditions
and finite prefixes, such that
\[
\widetilde A(i) \le C\,\rho_5^i \quad \text{for all } i.
\]
Dividing by $2^{B(i)}$ gives
\[
\frac{\widetilde A(i)}{2^{B(i)}} \le \frac{C\,\rho_5^i}{3^i} \longrightarrow 0,
\]
which completes the proof.
\end{proof}

Thus the exponential growth of admissible division words is strictly dominated
by the exponential growth of the denominator $2^{B(i)} \sim 3^i$.

% ============================================================
\section{Affine Constraint Fibers}

\subsection{Congruence Stability of Division Words}

\begin{theorem}[Congruence Stability of Division Words]
Let $\omega$ be admissible and let $r_\omega$ satisfy
\[
r_\omega \equiv -c(\omega)\cdot (3^m)^{-1} \pmod{2^{D(\omega)}}.
\]
Then for all $x \equiv r_\omega \pmod{2^{D(\omega)}}$, the first $m$
iterations of the accelerated Collatz map realize exactly the division word $\omega$.
\end{theorem}

\subsection{Residue Class Representation of Division Words}

\begin{theorem}[Affine Constraint Representation]
Let $\omega$ be an admissible division word of length $m$.

Then there exists a unique residue $r_\omega \pmod{2^{D(\omega)}}$ such that
\[
x \text{ realizes } \omega \iff x \equiv r_\omega \pmod{2^{D(\omega)}}.
\]

Thus admissible division words correspond to affine congruence constraints
of modulus $2^{D(\omega)}$.
\end{theorem}

\begin{proof}
From the affine representation,
\[
F_\omega(x) = \frac{3^m x + c(\omega)}{2^{D(\omega)}},
\]
integrality is equivalent to
\[
3^m x + c(\omega) \equiv 0 \pmod{2^{D(\omega)}}.
\]

Since $3^m$ is odd, it is invertible modulo $2^{D(\omega)}$, yielding
\[
x \equiv -c(\omega) \cdot (3^m)^{-1} \pmod{2^{D(\omega)}}.
\]

This determines a unique residue modulo $2^{D(\omega)}$.

Conversely, any $x$ in this class satisfies the integrality condition,
and therefore realizes the division word $\omega$.

Thus admissible division words are in one-to-one correspondence with affine
congruence constraints modulo powers of two.
\end{proof}

% ============================================================
\section{Exponential Sparsification}

\begin{lemma}[Constraint Refinement]
Each admissible division word $\omega$ imposes a congruence constraint
modulo $2^{D(\omega)}$, where $D(\omega) \ge A020914(m)$.
\end{lemma}

\begin{lemma}[Subcritical Symbolic Growth]
The number of admissible division words of length $m$ satisfies
\[
\#\{\omega : |\omega|=m\} = O(\rho^m)
\quad \text{for some } \rho < 3.
\]
\end{lemma}

\begin{proof}
This follows from Theorem~\ref{thm:subcritical-growth}, which bounds the growth
of admissible configurations generated by the increment grammar.
\end{proof}

\begin{theorem}[Exponential Sparsification]
The set of integers admitting a division word of length $m$ occupies at most
a proportion on the order of
\[
\frac{\#\{\text{admissible words of length } m\}}{2^{A020914(m)}}.
\]
In particular, this proportion decays exponentially in $m$.
\end{theorem}

\begin{lemma}[Exponential Tail Decay]
Let $E_m$ denote the set of integers that do not admit a descent prefix
of length at most $m$.

Then $\mu(E_m)$ decays exponentially in $m$.
\end{lemma}

\begin{proof}
An integer fails to admit a descent prefix of length $m$ only if it avoids
all admissible convergent division words of length at most $m$.

Each such word imposes a congruence constraint modulo $2^{D(\omega)}$,
with $D(\omega) \ge A020914(m)$.

By subcritical symbolic growth, the number of such constraints grows
strictly slower than $2^{A020914(m)}$, so the total proportion of integers
covered by these constraints approaches $1$ exponentially fast.

Hence the complement $E_m$ decays exponentially.
\end{proof}

% ============================================================
\section{Density-One Descent}

\begin{lemma}[Exhaustion by Descent Prefixes]
Let $S_m$ denote the set of integers that admit no convergent division word
of length $\le m$. Then
\[
S_{m+1} \subseteq S_m,
\quad\text{and}\quad
\bigcap_{m=1}^\infty S_m
\]
is precisely the set of integers that admit no finite descent prefix.
\end{lemma}

\begin{theorem}[Density-One Descent]
The set of integers that do not admit any finite convergent division word
has natural density zero.
\end{theorem}

\begin{proof}[Proof of Theorem]
For each $m$, the complement of $S_m$ is the union of all residue classes
corresponding to convergent division words of length $\le m$.

Each convergent division word of length $m$ defines a residue class
of density $2^{-D(\omega)} \le 2^{-B(m)}$, and these classes are disjoint.

By the exponential sparsification result, the total number of such words
grows like $\widetilde A(m) \le C \rho^m$ with $\rho < 3$, while
$2^{B(m)} \sim 3^m$. Hence the total density contributed by words of length
$m$ decays exponentially.

Summing over all lengths $\le m$, the total density of integers admitting
a convergent division word of length at most $m$ tends to $1$ as $m \to \infty$.
Equivalently, $\mu(S_m) \to 0$.

Therefore
\[
\mu\!\left(\bigcap_{m=1}^\infty S_m\right) = 0,
\]
which proves the theorem.
\end{proof}

\begin{corollary}
The set of integers whose trajectory under $f$ eventually decreases
has natural density one.
\end{corollary}

\begin{remark}
Density emerges as a consequence of affine growth imbalance.
\end{remark}

% ============================================================
\section{Discussion and Outlook}

% Interpretation of result
% Relation to classical approaches
% Remaining gap (zero-density exceptional set)
% Possible extensions

% ============================================================
\clearpage
\bibliographystyle{unsrtnat}
\bibliography{src/latex/references}

\end{document}